\section{The actual current embedding dimension}
\label{sec:hironaka-current-embedding-dimension}

The rank of the complete reduced current ideal detects singular points
even when that ideal has order one. The following construction extends
the exact Jacobian criterion of \cite{demailly-cadg}, II.4.31, to every
embedding-dimension superlevel and to its actual global maximum locus.
The finite-minor presentation agrees with the support description in
\cite{bierstone-milman97-invent-math}, Proposition~10.1; the present
construction uses the entire original reduced ideal at every nearby stalk.
It is a first geometric layer for the current-centre condition in
\cite{bierstone-milman97-invent-math}, Theorem~1.6(1). It does not assert
that this first value decreases under arbitrary singular-centre blowups;
normal flatness and the higher Hilbert--Samuel entries remain necessary.

\begin{definition}[Embedding dimension of the complete original ideal]
\label{hir:current-embedding-dimension}
\lean{Hironaka.HironakaCurrentEmbeddingDimension.embeddingDimension}
\leanok
\uses{rem:tangent-space}
For a set $A$ in an $n$-dimensional complex manifold and $x$ in the
ambient manifold, let $e(A,x)$ be the dimension of the common kernel,
in the actual preferred chart at $x$, of the differentials of all
holomorphic germs in the complete vanishing ideal $I_{A,x}$.
\end{definition}

\begin{lemma}[An actual chart-independent finite value]
\label{hir:embedding-dimension-chart-independent}
\lean{Hironaka.HironakaCurrentEmbeddingDimension.embeddingDimension_eq_finrank_tangentSpaceIn,Hironaka.HironakaCurrentEmbeddingDimension.embeddingDimension_le,Hironaka.HironakaCurrentEmbeddingDimension.embeddingDimension_eq_of_regularAtDim}
\leanok
\uses{hir:current-embedding-dimension,rem:tangent-space}
The same value is the dimension of that actual common kernel in any
chart containing $x$. It lies between $0$ and $n$. At a genuine regular
point of dimension $d$ it equals $d$.
\end{lemma}
\begin{proof}
The actual chart-transition derivative is an invertible linear map.
The chain rule carries the common kernel for all original ideal germs
to the common kernel in the other chart. Restricting this invertible derivative gives a linear isomorphism
between the two common kernels. This proves equality of their dimensions. The common kernel
is a subspace of $\mathbb C^n$, so its dimension is at most $n$.
At a regular point the actual flattening identifies it with the image
of the differential of the dimension-$d$ parametrization. That
differential is injective and its range has dimension $d$.
\end{proof}

\begin{lemma}[The actual complete-ideal rank dictionary]
\label{hir:embedding-dimension-cartan-rank-nullity}
\uses{hir:embedding-dimension-chart-independent,ana:stokes-cartan-generators,rem:tangent-space}
\lean{Hironaka.HironakaCurrentEmbeddingDimension.rank_add_embeddingDimension_of_cartanGenerators}
\leanok
Let holomorphic functions $s_1,\ldots,s_N$ generate the entire ideal of $A$
at every point of an open neighbourhood $U$ and cut out $A$ there. At
$x\in A\cap U$, in every genuine chart containing $x$,
\[
 \dim\operatorname{span}\{ds_i(x):1\le i\le N\}+e(A,x)=n.
\]
\end{lemma}
\begin{proof}
Each $s_i$ is an original ideal germ, so a vector in the kernel of every
original ideal differential lies in every $\ker ds_i(x)$. Conversely, an
arbitrary ideal germ has a finite holomorphic expression $\sum_i a_i s_i$
near $x$. Since all $s_i(x)=0$, its differential is $\sum_i a_i(x)ds_i(x)$.
Thus the common kernel of all actual ideal differentials is exactly the
common kernel of these finitely many $ds_i(x)$. Chart independence identifies
its dimension with $e(A,x)$. Rank-nullity for the linear map having these
covectors as its rows gives the displayed equality: its rank equals the
dimension of their span.
\end{proof}

\begin{lemma}[Every actual embedding-dimension superlevel is analytic]
\label{hir:embedding-dimension-superlevels-analytic}
\lean{Hironaka.HironakaCurrentEmbeddingDimension.isComplexAnalyticSubset_embeddingDimension_superlevel}
\leanok
\uses{hir:embedding-dimension-chart-independent,hir:embedding-dimension-cartan-rank-nullity,ana:stokes-cartan-generators,rem:tangent-space,ana:stokes-singular-locus-analytic}
Let $A$ be an arbitrary analytic set. For every integer $q\ge0$,
the actual whole locus $\{x\in A:e(A,x)\ge q\}$ is analytic, without
pure dimensionality, a supplied equation family or a constant-rank
certificate.
\end{lemma}
\begin{proof}
If $q>n$ the actual subspace dimension bound makes the locus empty.
Otherwise, at an arbitrary ambient point, produce actual Cartan
generators $s_1,\ldots,s_N$ of the entire original ideal on one open
neighbourhood. They cut out $A$ on that neighbourhood. At a point of
$A$, differentiating an arbitrary holomorphic ideal combination gives
a linear combination of the $ds_i$, since all $s_i$ vanish at that
point. Thus the actual common kernel of all ideal differentials is
exactly the common kernel of these finitely many differentials.
Rank-nullity gives
\[
 e(A,x)+\dim\operatorname{span}\{ds_i(x)\}=n.
\]
Consequently $e(A,x)\ge q$ if and only if the latter rank is less than
$n-q+1$, equivalently all $(n-q+1)$-minors of the actual Jacobian
vanish. In the true chart this whole locus is therefore cut out by
the original generators and the finite family of these minors.
Their derivatives and determinants are holomorphic on the same
neighbourhood, so this is a genuine local analytic presentation.
The chart-independent dimension equality proves it represents the
whole stated intrinsic locus. For $q=0$ the excessive-size minors are
automatically zero and the locus is precisely $A$; empty generator
families and $n=0$ are included in the same finite presentation.
\end{proof}

\begin{lemma}[The entire pure singular locus is the first strict superlevel]
\label{hir:pure-singular-embedding-dimension-threshold}
\lean{Hironaka.HironakaCurrentEmbeddingDimension.regularAtDim_iff_embeddingDimension_le,Hironaka.HironakaCurrentEmbeddingDimension.diff_regularLocus_eq_embeddingDimension_superlevel}
\leanok
\uses{hir:embedding-dimension-superlevels-analytic,hir:embedding-dimension-cartan-rank-nullity,rem:pure-singular-locus,hir:embedding-dimension-chart-independent}
\source{demailly-cadg:II.4.31}
Let $A$ be analytic and pure of dimension $d\le n$. At each $x\in A$,
regularity of dimension $d$ is equivalent to $e(A,x)\le d$.
Hence $A\setminus\operatorname{Reg}_d A$ is exactly
$\{x\in A:e(A,x)\ge d+1\}$.
\end{lemma}
\begin{proof}
At a regular-$d$ point the preceding actual tangent-space theorem gives
$e(A,x)=d$. Conversely if $e(A,x)\le d$, choose actual all-stalk Cartan
generators at $x$. The true kernel dictionary and rank-nullity give
rank at least $n-d$ for their differentials, hence a nonzero
$(n-d)$-minor. The actual pure-dimensional Jacobian criterion then
produces a flattening of the entire set $A$ at this same point of
dimension $d$. It uses all original generators to identify the
selected equations' smooth zero set with $A$, by regular density and
the holomorphic identity theorem. Negating this exact equivalence
gives the displayed whole singular-locus equality.
\end{proof}

\begin{theorem}[An internally produced singular maximum locus]
\label{hir:actual-embedding-dimension-maximum-singular-locus}
\lean{Hironaka.HironakaCurrentEmbeddingDimension.exists_maximum_embeddingDimension_singular_locus}
\leanok
\uses{hir:embedding-dimension-superlevels-analytic,hir:pure-singular-embedding-dimension-threshold}
If a pure analytic set $A$ of dimension $d\le n$ has a singular point,
its actual embedding dimension has an internally attained global
maximum $q>d$. Its entire maximum locus is nonempty, analytic, and
contained in the actual current singular locus of $A$. No compactness,
attained-maximum or selected-component hypothesis is needed.
\end{theorem}
\begin{proof}
The actual values on $A$ lie in the finite set $\{0,\ldots,n\}$.
A singular point makes their image nonempty, so this finite image has
a maximum $q$ attained at an actual point of $A$. The singular point's
value is at least $d+1$, hence $q>d$. Since every actual value on $A$
is at most $q$, the maximum locus is exactly the internally produced
superlevel $\{x\in A:e(A,x)\ge q\}$. Its analyticity is the preceding
genuine finite-minor producer. At each point of the locus $q>d$, so
the exact singular-threshold theorem places that point in
$A\setminus\operatorname{Reg}_d A$. This is a produced analytic
maximum locus; it does not assume that the locus is smooth or claim
a resolution/descent theorem from this first value alone.
\end{proof}
